\begin{lemma}[Reconstruction after a component outside the penultimate prefix]
Fix a minimum hierarchy counterexample and a raw
\noderef{HierarchySituationTwoCaseThreeData} \(D\).  Put
\[
 B=T(y_{p-2})=\noderef{HierarchyFullTailPrefix}(C,\ell-2),
 \qquad
 B^+=T(y_{p-1})=\noderef{HierarchyFullTailPrefix}(C,\ell-1),
\]
and let \(y_p\) be the terminal vertex stored by \(D\).  Suppose that a
color \(\beta\) is missing on \(B\), is absent from every exact edge of
the full tail above \(T_{n,q}^{*}\), and satisfies one of the following
three alternatives:
\begin{enumerate}
\item \(q>0\) and
      \(\beta\notin D_{n,q}\cup\Gamma^q\);
\item \(q=0\) and
      \(\beta\notin D_n\cup\Gamma^0\), where \(D_n\) is the raw terminal
      ladder stored set;
\item \(\beta\in D_{n,q}\) and
      \(\Gamma^q_\beta\) is disjoint from the colors of all exact edges
      of the full tail.
\end{enumerate}
Assume that \(\beta\) is not missing at \(y_p\).  Let
\[
 Q=P_{y_p}(\alpha,\beta,\varphi_n),
 \qquad
 \sigma_n=\varphi_n/Q,
\]
where \(\alpha=D.\alpha\), and assume that every vertex of \(Q\) is
outside \(B^+\).  Then \(Q\) contains \(y_p\), so it is not outside the
full tree \(T\).  Nevertheless there is a
\noderef{HierarchyReplacementPresentation} \(P\) on the same exact
ordered tree \(T\), with \(P.\mathit{coloring}=\sigma_n\), and a
\noderef{HierarchySituationTwoCaseThreeNormalizedData} \(N\) for \(P\)
whose normalized color is \(N.\alpha=\beta\).  In particular the returned
presentation contains an actual corresponding same-rung ETT, maximum
property, a rebuilt good hierarchy and counterexample, and the transported
minimum context with all six lower-rung assertions.
\end{lemma}
\begin{proof}
Write \(H=C.\mathit{hierarchy}\), let \(q\) be its last divider index,
put \(A=T_{n,q}^{*}\), and retain the exact displayed indices from \(D\).
The two position equations in
\noderef{HierarchySituationTwoCommonTailData} identify \(B^+\) with the
prefix obtained by deleting precisely the last edge \(e_p\) and last
vertex \(y_p\) from \(T\).  They also put \(y_p\) in
\(\noderef{TreeSequenceVertexSet}(T)\).  The root of a
\noderef{TwoColorComponent} belongs to its own component.  Hence
\(y_p\in Q\), proving both the first asserted membership and the fact that
\(Q\) meets \(T\).  The assumed disjointness says only that \(Q\) is
outside \(B^+\), as required; it does not make \(Q\) an outside component
of \(T\).

The raw datum says that \(\alpha\) is missing at \(y_p\), whereas the
present hypothesis says that \(\beta\) is not missing there.  Thus
\(\alpha\ne\beta\).  Use the singleton root set \(\{y_p\}\) in the
definition of \noderef{OutsideKempeRecoloring}.  The component-disjointness
hypothesis gives its outside clause, and the pointwise equation defining
\noderef{TwoColorComponentSwap} gives its coloring clause.  Therefore
\[
 \sigma_n\text{ is an outside Kempe recoloring of }B^+,
                                                               \tag{1}
\]
but, by the preceding paragraph, (1) is deliberately not asserted for
\(T\).

We next repackage the old hierarchy on the exact prefix \(B^+\).  The
common-tail inequality \(\ell\ge2\), together with the definition of
\noderef{HierarchyFullTailPrefix}, gives
\[
 A\prec B^+\prec T,
 \qquad |T|=|B^+|+1.                                      \tag{2}
\]
Apply \noderef{GoodHierarchyTerminalPrefixRepackage} to (2).  It gives a
positive-rung ETT \(\mathcal E^+\) on \(B^+\), with maximum property, the
same certified ladder series and rung number as the original ETT, a base
with the same core and PE closure, and a good hierarchy \(H^+\).  At each
divider its ordinary and starred stages, computed stored set, stage-missing
set, and Gamma pair are the corresponding old objects; the segment-color
sets are the exact restrictions to \(B^+\).

Now apply \noderef{HierarchyLEMStable} to \(\mathcal E^+\), \(H^+\), and
the outside recoloring (1).  Its lower-rung premise is exactly the
lower-rung field of the supplied
\noderef{HierarchyMinimumCounterexampleContext}; positivity and maximum
property came from the prefix repackaging.  The result yields
\[
 \begin{split}
 &\sigma_n\text{ proper and exactly-one-uncolored},\\
 &\sigma_n\text{ stable on the terminal ladder tree and invariant on }B^+,
                                                               \tag{3}\\
 &\text{a corresponding maximum ETT and the copied good hierarchy on }B^+.
 \end{split}
\]
This is the only invocation of the outside-Kempe hierarchy lemma; its tree
argument is \(B^+\), never \(T\).

We transport the old terminal augmentation route, including its last
individual edge.  Let \(T_n\) be the last ladder tree.  In the PE case the
terminal constructor in \noderef{TashkinovSeriesStep}, unfolded through
\noderef{TashkinovClosure}, supplies a Tashkinov-augmentation route from
\(T_n\) to the final closure.  In the SE or RE case it supplies the exact
singleton connecting edge \(f_n\), the exact seed \(T_n+f_n\), and a route
from that seed to the final closure.  Since the ETT places \(T\) as an
exact initial segment of that closure,
\noderef{TashkinovAugmentationSequenceRestrict} truncates the route at
the exact \(T\).  It also truncates it at \(B^+\).  Thus the final step of
the retained route is the displayed pair \((e_p,y_p)\), not an isomorphic
or parallel replacement edge.

The invariant part of (3) gives equality of the two missing-color sets at
every vertex of \(B^+\).  Every route edge preceding \(e_p\) is incident
with \(B^+\), and hence is fixed by (1).  For \(e_p\), the raw terminal
missingness of \(\alpha\) and the incidence of \(e_p\) with \(y_p\) show
\(\varphi_n(e_p)\ne\alpha\).  The candidate \(\beta\) is absent from every
full-tail edge, so in particular
\(\varphi_n(e_p)\ne\beta\).  Consequently the component swap fixes the
individual edge \(e_p\) as well.  We may therefore apply
\noderef{TashkinovAugmentationSequenceColoringTransport}, with \(B^+\) as
the penultimate prefix and \(T\) as the target.  Its two pointwise premises
are exactly the missing-set equality from (3) and the edge equalities just
proved.  The result is the identical terminal augmentation route under
\(\sigma_n\), including \(e_p\); call this conclusion (4).

The stable ladder-tree clause of (3), together with properness and exact
one-uncoloredness, may now be passed to \noderef{Td}.  We use the original
ETT, whose terminal ladder data agree with those of \(\mathcal E^+\), and
again use the lower-rung field of the minimum context.  This gives
\[
                 \noderef{ModColoring}(\mathcal E,\sigma_n).   \tag{5}
\]
For PE, (4) is the first eligibility branch of the universal clause in
(5).  For SE or RE, the retained exact singleton \(f_n\), its exact
one-edge seed, and (4) give the second branch.  Hence (5) produces an
actual \noderef{CorrespondingExtendedTashkinovTree} on the full exact
\(T\).  Applying \noderef{MaximumPropertyTransport} to this certificate
and the old maximum property gives maximum property for the replacement
ETT.

It remains to rebuild the hierarchy and the minimum presentation, rather
than merely naming alternative (a).  Copy the old core, closure, divider
count, exact stages, and divider positions to the replacement ETT.  All
nonterminal ordinary stages and every starred stage lie in \(B^+\), so (3) identifies
their vertexwise missing sets.  The component swap changes no displayed
edge of \(T\): the prefix edges are fixed by (1), while \(e_p\) was checked
separately above.  Thus every computed stored set is unchanged.  For each
stored color potentially affected by the swap, its first missing occurrence
is already in $B^+$: if it is $\alpha$, use the raw earlier witness in $B$,
and if it is $\beta$, use the candidate witness in $B$.  Every other stored
color is untouched by the two-color swap.
Consequently every hierarchy segment-color set is unchanged.  Copying the
old Gamma function now verifies, one field at a time, the pair cardinality,
avoidance, pairwise disjointness, unique refresh, PE-base clause, and
minus-closed clauses of \noderef{GoodHierarchy}.  This constructs a
same-good replacement hierarchy on \(T\), without applying
\noderef{HierarchyLEMStable} to \(T\).

The candidate \(\beta\) is missing at a vertex of \(B\).  That vertex lies
in \(B^+\), so (3) preserves its missingness.  Swapping the component at
\(y_p\), where \(\alpha\) was missing and \(\beta\) was present, makes
\(\beta\) missing at \(y_p\).  These are distinct displayed vertices by
(2); hence the replacement terminal vertex set is not elementary.  With
the maximum and good-hierarchy data already established, this gives the
replacement \noderef{HierarchyCounterexample}.

Finally transport the minimum-context fields explicitly.  Properness and
exact one-uncoloredness are in (3).  The six lower-rung assertions and the
fewer-divider induction are graph- and rung-local universal hypotheses, so
copy verbatim.  Same-goodness identifies the divider count and every exact
stage.  It therefore identifies both coordinates of
\noderef{HierarchyCounterexampleMeasure}; the old lexicographic-minimality
field proves the new one after this equality.  Stability on \(T_n\) in (3)
identifies its tree-missing set, so the old \(2n+11\) estimate is retained.
These are precisely all fields of
\noderef{HierarchyMinimumCounterexampleContext}, and hence all fields of
the required \noderef{HierarchyReplacementPresentation}.

For the normalized datum, retain all exact indices and the same last-edge
color \(\theta\), since the individual edge \(e_p\) is fixed.  Choose as
the earlier witness a displayed vertex of \(B\) at which \(\beta\) is
missing; its index is strictly before that of \(y_{p-1}\) by (2).  The two
missing-vertex assertions were proved above.  Missingness on (B), full
tail nonuse, and the three-way normalized alternative are respectively
the three hypotheses on \(\beta\), transported through the established
prefix invariance and edge equality.  This fills every field of
\noderef{HierarchySituationTwoCaseThreeNormalizedData}, with normalized
color \(\beta\), and completes the reconstruction.
\end{proof}
